\exercisesection{\S~4}

\begin{enumerate}
\def\labelenumi{\arabic{enumi})}
\tightlist
\item
  取 G 为群 \(\mathbf{SU}(n,\mathbf{C})\) 或 \(\mathbf{U}(n,\mathbf{H})\) 之一；把 \(g_{C}\) 相应地与 \(\mathfrak{sl}(n,\mathbf{C})\) 或 \(\mathfrak{sp}(2n,\mathbf{C})\) 等同，参见~§3，节 4 及习题 3。我们采用第 VIII 章，§13，节 1 与节 3 的记号，其中 k=C。
\end{enumerate}

\begin{enumerate}
\def\labelenumi{\alph{enumi})}
\item
  证明 G 中由复系数对角矩阵组成的子群 T 是一个极大环面，且 \(\mathrm{L}(\mathrm{T})_{(\mathbf{C})} = \mathfrak{h}\)。
\item
  把 X(T) 与 \(\mathfrak{h}^*\) 的一个子群等同起来（借助同态 \(\delta\)）。证明线性形式 \(\varepsilon_i\) 与支配权 \(\varpi_j\) 都属于 X(T)。若 \(t = \mathrm{diag}(t_1, \ldots, t_n) \in \mathrm{T}\)，则 \(\varepsilon_i(t) = t_i\) 且 \(\varpi_j(t) = t_1 \ldots t_j\)（对 \(1 \leq i, j \leq n\)）。
\item
  由 b) 给出群 \(\mathbf{SU}(n,\mathbf{C})\) 与 \(\mathbf{U}(n,\mathbf{H})\) 单连通（参见~§3，习题 4 与 6）这一事实的另一证明。
\end{enumerate}

\begin{enumerate}
\def\labelenumi{\arabic{enumi})}
\setcounter{enumi}{1}
\tightlist
\item
  取 \(\mathrm{G} = \mathbf{SO}(n, \mathbf{R})\)，其中 \(n \geq 3\)；当 n 为奇数时置 \(n = 2l + 1\)，当 n 为偶数时置 n = 2l。代数 \(g_{C}\) 可与 \(\mathfrak{o}(n, \mathbf{C})\) 等同；我们采用第 VIII 章，§13，节 2 与节 4 的记号。用 \((f_i)_{1 \leq i \leq n}\) 表示 \(R^n\) 的典则基。置 \(e_j = \frac{1}{\sqrt{2}}(f_{2j-1} + if_{2j})\) 与 \(e_{-j} = \frac{1}{\sqrt{2}}(f_{2j-1} - if_{2j})\)（对 \(1 \leq j \leq l\)），当 n 为奇数时再置 \(e_0 = i\sqrt{2}f_{2l+1}\)；当 n 为偶数（相应地奇数）时，选取 \(C^n\) 的由 \(e_1, \ldots, e_l, e_{-l}, \ldots, e_{-1}\)（相应地 \(e_1, \ldots, e_l, e_0, e_{-l}, \ldots, e_{-1}\)）给出的 Witt 基。
\end{enumerate}

\begin{enumerate}
\def\labelenumi{\alph{enumi})}
\tightlist
\item
  设 \(H_{i}\) 是 \(R^{n}\) 中由 \(f_{2i-1}\) 与 \(f_{2i}\) 生成的子空间（\(1 \leq i \leq l\)）；证明 G 中由满足 \(g(\mathrm{H}_{i}) \subset \mathrm{H}_{i}\) 且 \(\det(g|\mathrm{H}_{i}) = 1\)（对 \(1 \leq i \leq l\)）的元素 g 组成的子群是 G 的一个极大环面 T，且 \(\mathrm{L}(\mathrm{T})_{(\mathbf{C})} = \mathfrak{h}\)。
\item
  把 X(T) 与 \(h^{*}\) 的一个子群等同起来（借助 \(\delta\)）。证明线性形式 \(\varepsilon_{i}\) 属于 X(T)；若 \(t \in T\) 且 t 在 \(H_{j}\) 上的限制是转角为 \(\theta_{j}\) 的旋转，则 \(\varepsilon_{j}(t) = e^{i\theta_{j}}\)。诸权 \(\varpi_{1}, \ldots, \varpi_{l-2}\)；\(2\varpi_{l-1}, 2\varpi_{l}, \varpi_{l-1} \pm \varpi_{l}\) 都属于 X(T)。若 n 为奇数，则 \(\varpi_{l-1}\) 属于 X(T)。
\item
  设 \(\tilde{\mathbf{G}} = \mathbf{Spin}(n, \mathbf{R})\)，并设 \(\varphi : \tilde{G} \to G\) 是典则覆盖。对 \(\boldsymbol{\theta} = (\theta_{1}, \ldots, \theta_{l}) \in \mathbf{R}^{l}\)，置 \(t(\boldsymbol{\theta}) = \prod_{i=1}^{l} (\cos\theta_{i} - f_{2i-1}f_{2i}\sin\theta_{i}) \in \tilde{\mathbf{G}}\)。证明诸 \(t(\boldsymbol{\theta})\)（\(\theta \in R^{l}\)）组成的集合是一个极大环面 \(\tilde{T}\)（\(\tilde{G}\) 中的），且 \(\varphi(\tilde{T}) = T\)。若把 X( \(\tilde{T}\) ) 与 \(h^{*}\) 的一个子群等同起来，则有 \(\varepsilon_{j}(t(\boldsymbol{\theta})) = e^{2i\theta_{j}}\)。
\item
  证明权 \(\varpi_{l-1}\) 与 \(\varpi_{l}\) 属于 X( \(\tilde{T}\) )；由此推出 \(\mathbf{Spin}(n,\mathbf{R})\) 是单连通的（参见~§3，习题 5）。
\end{enumerate}

\begin{enumerate}
\def\labelenumi{\arabic{enumi})}
\setcounter{enumi}{2}
\tightlist
\item
  \begin{enumerate}
  \def\labelenumii{\alph{enumii})}
  \tightlist
  \item
    证明 \(\sigma: \mathrm{A} \mapsto \overline{\mathrm{A}}\) 这一 \(\mathbf{SU}(n, \mathbf{C})\) 的自同构当 \(n \geq 3\) 时不是内自同构。\(\mathbf{SU}(n, \mathbf{C})\) 的每个非内自同构都具有 \((\operatorname{Int} g) \circ \sigma\)（\(g \in \mathbf{SU}(n, \mathbf{C})\)）的形式。
  \end{enumerate}
\end{enumerate}

\begin{enumerate}
\def\labelenumi{\alph{enumi})}
\setcounter{enumi}{1}
\tightlist
\item
  证明对每个 \(A_{n}\) 型（\(n \geq 2\)）的概单连通紧群 G，群 \(\operatorname{Aut}(\mathrm{G})/\operatorname{Int}(\mathrm{G})\) 都是 2 阶循环群（参见~§3，习题 4）。
\end{enumerate}

\begin{enumerate}
\def\labelenumi{\arabic{enumi})}
\setcounter{enumi}{3}
\tightlist
\item
  设 \(n\) 为整数 \(\geq 2\)。
\end{enumerate}

\begin{enumerate}
\def\labelenumi{\alph{enumi})}
\item
  设 \(g \in \mathbf{O}(2n, \mathbf{R})\) 且 \(\det g = -1\)；证明 \(\operatorname{Int} g\) 这一 \(\mathbf{O}(2n, \mathbf{R})\) 的自同构诱导出 \(\mathbf{SO}(2n, \mathbf{R})\) 的一个非内自同构。
\item
  当 \(n \geq 2\) 时，群 \(\operatorname{Aut}(\mathbf{SO}(2n, \mathbf{R}))\) 等于 \(\operatorname{Int}(\mathbf{O}(2n, \mathbf{R}))\)（它同构于 \(\mathbf{O}(2n, \mathbf{R}) / \{\pm I_{2n}\}\)）。
\item
  对 \(\mathbf{Spin}(2n,\mathbf{R})\) 与 \(\mathbf{SO}(2n,\mathbf{R})/\{\pm I_{2n}\}\) 建立类似的结果。若 n 为偶数且 \(\neq 2\)，则群 \(\mathbf{Spin}(2n,\mathbf{R})/\{1,\varepsilon\}\)（§3，习题 5）的每个自同构都是内自同构。
\end{enumerate}

\begin{enumerate}
\def\labelenumi{\arabic{enumi})}
\setcounter{enumi}{4}
\tightlist
\item
  设 \(\mathbf{R}\) 是一个既约且不可约的根系。
\end{enumerate}

\begin{enumerate}
\def\labelenumi{\alph{enumi})}
\item
  证明 R 中长度最大的根组成的集合是 R 的一个对称闭子集，且在 W(R) 下稳定。
\item
  证明 R 的每个非空、闭、对称且在 W(R) 下稳定的子集 P 都等于 R，或等于 R 中长度最大的根组成的集合。
\item
  设 \(P \neq R\)。证明根系 P 为 \(D_{l}\) 型（相应地 \((\mathrm{A}_{1})^{l}, \mathrm{D}_{4}, \mathrm{A}_{2}\)）当 R 为 \(B_{l}\) 型（相应地 \(C_{l}, F_{4}, G_{2}\)）时。
\end{enumerate}

\begin{enumerate}
\def\labelenumi{\arabic{enumi})}
\setcounter{enumi}{5}
\tightlist
\item
  设 \(\mathrm{H}\) 是 \(\mathrm{G}\) 的包含 \(\mathrm{N}_{\mathrm{G}}(\mathrm{T})\) 的闭子群。
\end{enumerate}

\begin{enumerate}
\def\labelenumi{\alph{enumi})}
\item
  证明 H 等于它在 G 中的正规化子，且 \(H/H_{0}\) 同构于 \(W/W_{H_{0}}(T)\)。
\item
  证明 \(\mathrm{R}(\mathrm{H}_{0},\mathrm{T})\) 在 W 下稳定。反之，若 K 是包含 T 的连通闭子群，且 \(\mathrm{R}(\mathrm{K},\mathrm{T})\) 在 W 下稳定，则 K 的正规化子包含 T 的正规化子。
\end{enumerate}

\(c)\) 设 \(\mathbf{G}\) 是概单的。证明 \(\mathbf{H}\) 等于 \(\mathrm{N_G(T)}\)，或等于 \(\mathbf{G}\)，或者（在同构意义下）处于下列情形之一：

\(\alpha\)（相应地 \(\alpha'\)）\(\mathrm{G} = \mathbf{Spin}(2l + 1, \mathbf{R})\)（相应地 \(\mathrm{G} = \mathbf{SO}(2l + 1, \mathbf{R})\)），且 \(\mathrm{H}_0\) 是 \(\mathbf{R}^{2l + 1}\) 中某个非零向量的稳定子群，它同构于 \(\mathbf{Spin}(2l, \mathbf{R})\)（相应地 \(\mathbf{SO}(2l, \mathbf{R})\)）；

\(\beta)\) \(\mathrm{G} = \mathrm{U}(l,\mathbf{H})\)，且 \(\mathrm{H}_0\) 是由对角矩阵组成的子群 D；\(\beta^{\prime})\) \(\mathrm{G} = \mathrm{U}(l,\mathbf{H}) / \{\pm 1\}\)，且 \(\mathrm{H}_0 = \mathrm{D} / \{\pm 1\}\)；

\(\gamma)\) G 是 \(\mathrm{F}_4\) 型的，且 \(\mathrm{H}_0\) 同构于 Spin(8,R)；

\(\delta)\) G 是 \(\mathrm{G}_2\) 型的，且 \(\mathrm{H}_0\) 是 §3 习题 7 中定义的子群（它同构于 \(\mathbf{SU}(3,\mathbf{C})\)）。

\begin{enumerate}
\def\labelenumi{\arabic{enumi})}
\setcounter{enumi}{6}
\tightlist
\item
  设 \(\tau : \mathbf{G} \to \mathbf{GL}(\mathbf{V})\) 是 \(\mathbf{G}\) 在一个有限维实向量空间上的连续表示。假定对所有 \(\lambda \in \mathrm{X}(\mathrm{T})\)，都有 \(\dim_{\mathbf{C}} \hat{V}_{\lambda} \leq 1\)。
\end{enumerate}

\begin{enumerate}
\def\labelenumi{\alph{enumi})}
\item
  证明表示 \(\tau\) 是一个由两两不同构的不可约表示组成的有限族 \((\tau_{i})_{1\leq i\leq s}\) 的直和，且其交换子 \(K_{i}\ (1\leq i\leq s)\) 同构于 R 或 C。
\item
  在 V 上存在使 \(\tau(\mathrm{G})\) 的运算均为 C-线性的 C-向量空间结构，当且仅当 \(K_{1},\ldots,K_{s}\) 都同构于 C；此时这种结构共有 \(2^{s}\) 个。
\end{enumerate}

\begin{enumerate}
\def\labelenumi{\arabic{enumi})}
\setcounter{enumi}{7}
\tightlist
\item
  设 H 是 G 的一个具有极大秩的连通闭子群，h 是它的李代数，X 是流形 G/H，V 是 X 在对应于 H 的类的点处的切空间；把 V 与商 g/h 等同起来。
\end{enumerate}

\begin{enumerate}
\def\labelenumi{\alph{enumi})}
\tightlist
\item
  若 j 是 X 上的一个殆复结构（《微分流形与解析流形，结果》，8.8.3），用 \(V''(j)\) 表示复向量空间 \(C \otimes V\) 中由满足 \(j(u) = -iu\) 的元素 u 组成的子空间，并用 \(\mathfrak{q}(j)\) 表示 \(g_{C}\) 中作为 \(V''(j)\) 的逆像的子空间，使得当 V 赋以由 j 诱导的 C-向量空间结构时，典则映射 \(V \to \mathfrak{g}_{\mathbf{C}} / \mathfrak{q}(j)\) 是一个 C-线性同构。证明映射 \(j \mapsto \mathfrak{q}(j)\) 是从 X 上在 G 下不变的殆复结构的集合，到 \(g_{C}\) 中满足下列条件的复子空间 p 的集合的一个双射：
\end{enumerate}

\[
\mathfrak {p} + \bar {\mathfrak {p}} = \mathfrak {g} _ {\mathbf {C}}\tag{1}
\]

\[
\mathfrak {p} \cap \bar {\mathfrak {p}} = \mathfrak {h} _ {\mathbf {C}}\tag{2}
\]

\[
[ \mathfrak {h}, \mathfrak {p} ] \subset \mathfrak {p}.\tag{3}
\]

\begin{enumerate}
\def\labelenumi{\alph{enumi})}
\setcounter{enumi}{1}
\item
  X 上存在这样的结构，当且仅当 H 在 V 上的伴随表示的诸不可约子表示的交换域都同构于 C；此时这种结构共有 \(2^{s}\) 个，其中 s 是 V 的不可约子表示的个数（利用习题 7）。
\item
  设 j 是 X 上在 G 下不变的殆复结构，\(\mathfrak{p} = \mathfrak{q}(j)\) 是相联的子空间。则 j 可积（即对应于 X 上的一个复解析流形结构，参见~《微分流形与解析流形，结果》，8.8.5 至 8.8.8），当且仅当 p 满足条件
\end{enumerate}

\[
[ \mathfrak {p}, \mathfrak {p} ] \subset \mathfrak {p}.\tag{4}
\]

\begin{enumerate}
\def\labelenumi{\alph{enumi})}
\setcounter{enumi}{3}
\item
  X 上存在在 G 下不变的复结构（即可积的殆复结构），当且仅当 H 是 G 中某个环面的中心化子（证明上面的条件 (1) 至 (4) 蕴含 p 是 \(g_{C}\) 的一个抛物子代数（第 VIII 章，§3，节 5））；此时这些复结构与 \(g_{C}\) 中那些是 \(h_{C}\) 与其幂幺根基的直和的抛物子代数 p 一一对应（参见~第 VIII 章，§3，节 4）。
\item
  证明 G/T 上在 G 下不变的复结构恰有 \(\operatorname{Card}(W)\) 个；若 \(\sigma\) 与 \(\sigma'\) 是两个这样的结构，则存在唯一的元素 \(w \in W\)，使 w 在 G/T 上的典则作用（通过内自同构）把 \(\sigma\) 变为 \(\sigma'\)。若 w 是 W 的非单位元素，则对 G/T 上任何在 G 下不变的复结构，w 在 G/T 上的作用都不是 C-解析的。
\item
  确定 \(S^{2}\) 上在 \(\mathbf{SO}(3)\) 下不变的复结构。
\end{enumerate}

\(g)^{*}\) 采用习题 8 \(d\)) 的记号，设 \(\mathrm{G}_c\) 是 \(\mathrm{G}\) 的复化，\(\mathrm{P}\) 是 \(\mathrm{G}_c\) 中以 \(\mathfrak{p}\) 为李代数的复李子群。证明典则映射 \(\mathrm{G / H} \to \mathrm{G}_c / \mathrm{P}\) 是复解析流形的一个同构.*

\begin{enumerate}
\def\labelenumi{\arabic{enumi})}
\setcounter{enumi}{8}
\tightlist
\item
  设 H 是 G 的一个连通闭子群，它具有极大秩、异于 G，且对这些性质而言是极大的。用 Z 表示商群 \(\mathrm{C}(\mathrm{H})/\mathrm{C}(\mathrm{G})\)。
\end{enumerate}

\begin{enumerate}
\def\labelenumi{\alph{enumi})}
\item
  我们有 \(\dim Z \leq 1\)；若 \(\dim Z = 0\)，则 Z 的阶为 2、3 或 5（化归到 G 概单且中心平凡的情形；应用第 VI 章，§4，习题 4）。
\item
  假定 G 是概单的且 Z 的阶为 3 或 5；置 \(z = \text{Card}(Z)\) 与 \(\pi = \text{Card}(\pi_1(H))/\text{Card}(\pi_1(G))\)。证明我们必定处于下列七种情形之一：
\end{enumerate}

\begin{enumerate}
\def\labelenumi{(\roman{enumi})}
\item
  G 是 \(\mathrm{G}_2\) 型的，H 是 \(\mathrm{A}_2\) 型的，且 \(z = 3, \pi = 1\)；
\item
  G 是 \(F_{4}\) 型的，H 是 \(A_{2} \times A_{2}\) 型的，且 \(z = \pi = 3\)；
\item
  G 是 \(E_{6}\) 型的，H 是 \(A_{2} \times A_{2} \times A_{2}\) 型的，且 \(z = \pi = 3\)；
\item
  G 是 \(E_{7}\) 型的，H 是 \(A_{2} \times A_{5}\) 型的，且 \(z = \pi = 3\)；
\item
  G 是 \(E_{8}\) 型的，H 是 \(A_{8}\) 型的，且 \(z = \pi = 3\)；
\item
  G 是 \(E_{8}\) 型的，H 是 \(A_{2} \times E_{6}\) 型的，且 \(z = \pi = 3\)；
\item
  G 是 \(E_{8}\) 型的，H 是 \(A_{4} \times A_{4}\) 型的，且 \(z = \pi = 5\)。
\end{enumerate}

（利用同上以及第 VI 章的图版；为计算 \(\pi\)，注意若 \(f\) 与 \(f'\) 分别表示 G 与 H 的连通指数，则 \(z\pi f = f'\)。）

\begin{enumerate}
\def\labelenumi{\alph{enumi})}
\setcounter{enumi}{2}
\tightlist
\item
  在上述每种情形中，确定群 H。
\end{enumerate}

\begin{enumerate}
\def\labelenumi{\arabic{enumi})}
\setcounter{enumi}{9}
\tightlist
\item
  我们沿用上一习题的记号。
\end{enumerate}

\begin{enumerate}
\def\labelenumi{\alph{enumi})}
\item
  假定 \(\dim Z = 1\)。则 G/H 上在 G 下不变的复结构恰有两个；存在 G 的一个自同构，它使 H 稳定并互换这两个结构（利用习题 8）。
\item
  确定 \(\mathbf{P}_n(\mathbf{C})\) 上在 \(\mathbf{SU}(n + 1,\mathbf{C})\) 下不变的复结构。
\item
  假定 \(\dim Z = 0\) 且 \(\operatorname{Card}(Z) \neq 2\)。证明 G/H 上存在在 G 下不变的殆复结构（若 z 是 C(H) 中在 G 内不是中心的元素，则 Int z 诱导出 G/H 的一个奇数阶自同构（习题 9）；利用习题 8 b)）。证明 G/H 上不存在在 G 下不变的复结构（利用习题 8 d)）。
\end{enumerate}

\(d\) ) \(\mathbf{S}_6\) 上不存在在 \(\mathbf{SO}(7,\mathbf{R})\) 下不变的复结构（利用 §3 的习题 7）。

\begin{enumerate}
\def\labelenumi{\arabic{enumi})}
\setcounter{enumi}{10}
\tightlist
\item
  我们沿用习题 9 与 10 的记号。
\end{enumerate}

\begin{enumerate}
\def\labelenumi{\alph{enumi})}
\item
  齐性空间 G/H 是对称的（§1，习题 8），当且仅当 Z 的阶为 2 或维数为 1。此时对称空间 G/H 是不可约的（同上）。
\item
  假定 \(\dim Z = 1\)；用 \(X\) 表示复解析流形 \(G / H\)。证明此时我们处于下列情形之一：
\end{enumerate}

\begin{enumerate}
\def\labelenumi{(\roman{enumi})}
\item
  \(\mathbf{G}\) 是 \(\mathrm{A}_l\) 型的，\(\mathrm{D(H)}\) 是 \(\mathrm{A}_{p - 1}\times \mathrm{A}_{l - p}\) 型的；\(\mathbf{X}\) 同构于格拉斯曼流形 \(\mathbf{G}_p(\mathbf{C}^{l + 1})\)。
\item
  G 是 \(B_{l}\) 型的，D(H) 是 \(B_{l-1}\) 型的；X 同构于 \(\mathbf{P}_{2l}(\mathbf{C})\) 中由一个分离二次型的取零所定义的子流形（光滑射影二次曲面）。
\end{enumerate}

(ii') G 是 \(D_{l}\) 型的，\(D(H)\) 是 \(D_{l-1}\) 型的；X 同构于 \(\mathbf{P}_{2l-1}(\mathbf{C})\) 中的一个光滑射影二次曲面。

\begin{enumerate}
\def\labelenumi{(\roman{enumi})}
\setcounter{enumi}{2}
\item
  G 是 \(C_{l}\) 型的，D(H) 是 \(A_{l-1}\) 型的；X 同构于 \(\mathbf{G}_{l}(\mathbf{C}^{2l})\) 中由 \(C^{2l}\) 上通常的交错双线性型的极大迷向子空间组成的子流形。
\item
  G 是 \(D_{l}\) 型的，\(\mathrm{D}(\mathrm{H})\) 是 \(A_{l-1}\) 型的；X 同构于 \(\mathbf{G}_{l}(\mathbf{C}^{2l})\) 中由 \(C^{2l}\) 上通常的对称双线性型的极大迷向子空间组成的子流形。
\item
  G 是 \(E_{6}\) 型的，D(H) 是 \(D_{5}\) 型的。
\item
  G 是 \(E_{7}\) 型的，D(H) 是 \(E_{6}\) 型的。
\end{enumerate}

\begin{enumerate}
\def\labelenumi{\alph{enumi})}
\setcounter{enumi}{2}
\tightlist
\item
  假定 \(\operatorname{Card}(Z)=2\)；给出可能情形的列表。若 G 是 \(A_{l}\)、\(B_{l}\)、\(C_{l}\) 或 \(D_{l}\) 型的，证明实流形 G/H 同构于一个格拉斯曼流形 \(\mathrm{G}_{p}(\mathrm{K}^{q})\)，其中 K=R、C 或 H。
\end{enumerate}

¶ 12) 假定 G 是单连通的；对所有 \(\alpha \in \mathrm{R}(\mathrm{G},\mathrm{T})\)，置 \(t_{\alpha} = \exp (\frac{1}{2} K_{\alpha})\)。置 \(\mathrm{N} = \mathrm{N}_{\mathrm{G}}(\mathrm{T})\)，并用 \(\varphi : \mathrm{N} \to \mathrm{W}\) 表示典则映射。设 B 是 \(\mathrm{R}(\mathrm{G},\mathrm{T})\) 的一个基；选取元素 \(n_{\alpha}\)（在 \((\mathrm{N} \cap \mathrm{S}_{\alpha}) - (\mathrm{T} \cap \mathrm{S}_{\alpha})\) 中，对所有 \(\alpha \in \mathrm{B}\)）。

\begin{enumerate}
\def\labelenumi{\alph{enumi})}
\item
  证明 \(\varphi(n_{\alpha}) = s_{\alpha}\) 且 \(n_{\alpha}^{2} = t_{\alpha}\)，从而 \(n_{\alpha}^{4} = 1\)。
\item
  设 \(\alpha\) 与 \(\beta\) 是 B 的两个不同元素，并设 \(m_{\alpha \beta}\) 是 \(s_{\alpha}s_{\beta}\) 在 W 中的阶。证明
\end{enumerate}

\[
\begin{array}{r l} {n _ {\alpha} n _ {\beta} = n _ {\beta} n _ {\alpha}} & {\mathrm{if} m _ {\alpha \beta} = 2} \\ {n _ {\alpha} n _ {\beta} n _ {\alpha} = n _ {\beta} n _ {\alpha} n _ {\beta}} & {\mathrm{if} m _ {\alpha \beta} = 3} \\ {(n _ {\alpha} n _ {\beta}) ^ {2} = (n _ {\beta} n _ {\alpha}) ^ {2}} & {\mathrm{if} m _ {\alpha \beta} = 4} \\ {(n _ {\alpha} n _ {\beta}) ^ {3} = (n _ {\beta} n _ {\alpha}) ^ {3}} & {\mathrm{if} m _ {\alpha \beta} = 6} \end{array}
\]

（例如若 \(m_{\alpha\beta}=3\)，则 \((s_{\alpha}s_{\beta})s_{\alpha}(s_{\alpha}s_{\beta})^{-1}=s_{\beta}\) 且 \(s_{\alpha}s_{\beta}(\alpha)=\beta\)；证明 \(n_{\alpha}n_{\beta}n_{\alpha}n_{\beta}^{-1}n_{\alpha}^{-1}n_{\beta}^{-1}\) 属于 \(S_{\beta}\)，并由 \(S_{\alpha}\cap S_{\beta}\cap T=\{e\}\) 得出结论）。

\begin{enumerate}
\def\labelenumi{\alph{enumi})}
\setcounter{enumi}{2}
\item
  由 b) 推出，存在唯一的截面 \(\nu: W \to N\)（\(\varphi\) 的），使得 \(\nu(s_{\alpha}) = n_{\alpha}\)，且 \(\nu(ww') = \nu(w)\nu(w')\)（当 \(l(ww') = l(w) + l(w')\) 时）（其中 \(l(w)\) 表示 w 关于生成系 \((s_{\alpha})_{\alpha \in \mathbb{B}}\) 的长度；参见~第 VI 章，§1，节 5，命题 5）。置 \(n_w = \nu(w)\)。
\item
  设 \(W^{*}\) 是 N 中由诸 \(n_{\alpha}\) 生成的子群；证明 \(W^{*} \cap T\) 是子群 \(T_{2}\)（T 中由阶 \(\leq 2\) 的元素组成），且 W 可与 \(W^{*}/T_{2}\) 等同。
\item
  设 \(w \in \mathrm{W}\) 满足 \(w^2 = 1\)；证明 \(n_w^2 = \prod_{\alpha \in \mathrm{R}_w} t_\alpha\)，其中 \(\mathrm{R}_w\) 是正根 \(\alpha\)（满足 \(w(\alpha) < 0\)）组成的集合（写 \(w = s_{\alpha_1} \ldots s_{\alpha_r}\)，其中 \(r = l(w)\) 且 \(\alpha_i \in \mathrm{B}\)；应用 c) 以及第 VI 章，§1，节 6，命题 17 的推论 2）。
\end{enumerate}

\(f\) ) 假定 \(\mathrm{G}\) 是概单的。设 \(c\) 是 \(\mathrm{W}\) 中的一个 Coxeter 变换，\(h\) 是 \(\mathrm{W}\) 的 Coxeter 数（第 VI 章，§1，节 11）。证明 \(n_c^h = \prod_{\alpha \in \mathbb{R}_+}t_\alpha\)（利用 \(c\)）、\(e\)）以及第 V 章，§6 的习题 2）。

\begin{enumerate}
\def\labelenumi{\arabic{enumi})}
\setcounter{enumi}{12}
\tightlist
\item
  \begin{enumerate}
  \def\labelenumii{\alph{enumii})}
  \tightlist
  \item
    选取 R(G, T) 的一个基 B，并用 R+ 表示 R(G, T) 的正根集合。证明 \(z_{G} = \prod_{\alpha \in R_{+}} \exp(\frac{1}{2} K_{\alpha})\) 是 C(G) 的一个元素，它与 T 和 B 的选取无关；我们有 \(z_{G}^{2} = e\)。
  \end{enumerate}
\end{enumerate}

\begin{enumerate}
\def\labelenumi{\alph{enumi})}
\setcounter{enumi}{1}
\item
  设 H 是另一个连通紧李群。则 \(z_{\mathrm{G}\times\mathrm{H}} = (z_{\mathrm{G}}, z_{\mathrm{H}})\)；证明若 \(f : G \to H\) 是李群的满射态射，则 \(f(z_{\mathrm{G}}) = z_{\mathrm{H}}\)。
\item
  置 \(\mathrm{R} = \mathrm{R}(\mathrm{G},\mathrm{T})\)；假定 \(\mathbf{G}\) 是单连通的，从而 \(\mathrm{X(T)}\) 可与 \(\mathrm{P(R)}\) 等同。证明同态 \(\chi \mapsto \chi (z_{\mathrm{G}})\)
\end{enumerate}

（从 X(T) 到 \(\{1,-1\}\)）的核是第 VI 章，§1，习题 8 中定义的子群 \(P'(R)\)。

\begin{enumerate}
\def\labelenumi{\alph{enumi})}
\setcounter{enumi}{3}
\tightlist
\item
  若 \(\mathrm{G} = \mathbf{S}\mathbf{U}(n,\mathbf{C})\)，则 \(z_{\mathrm{G}} = (-1)^{n + 1}I_n\)；
\end{enumerate}

若 \(\mathrm{G} = \mathbf{S}\mathbf{U}(n,\mathbf{H})\)，则 \(z_{\mathrm{G}} = -I_n\)；

若 \(\mathrm{G} = \mathbf{Spin}(n, \mathbf{R})\) 且 \(n \equiv 3, 4, 5\) 或 6 (mod 8)，则 \(z_{G} = -1\)；

若 \(\mathrm{G} = \mathbf{Spin}(n, \mathbf{R})\) 且 \(n \equiv 0, 1, 2, 7 \pmod{8}\)，则 \(z_{G} = 1\)；

若 G 是 \(E_{6}\)、\(E_{8}\)、\(F_{4}\) 或 \(G_{2}\) 型的，则 \(z_{G} = e\)；

若 G 是 \(E_{7}\) 型单连通的，则 \(z_{G}\) 是 \(\mathrm{C}(\mathrm{G})\) 的唯一的非单位元素。

（利用第 VI 章，§4，习题 5。）

\begin{enumerate}
\def\labelenumi{\arabic{enumi})}
\setcounter{enumi}{13}
\tightlist
\item
  假定群 G 是概单的；用 h 表示 R(G,T) 的 Coxeter 数（第 VI 章，§1，节 11）。若存在 G 的一个极大环面 S，使 g 属于 \(\mathrm{N}_{\mathrm{G}}(\mathrm{S})\)，且 g 在 \(\mathrm{W}_{\mathrm{G}}(\mathrm{S})\) 中的类是一个 Coxeter 变换（同上），则称 G 的元素 g 为 Coxeter 元。
\end{enumerate}

\begin{enumerate}
\def\labelenumi{\alph{enumi})}
\item
  证明任意两个 Coxeter 元都共轭（按节 5 命题 10 的推论的证明方式论证）。
\item
  每个 Coxeter 元 g 都是正则的，且满足 \(g^{h} = z_{G}\)，其中 \(z_{G}\) 是习题 13 中定义的 C(G) 的元素；特别地，视 \(z_{G}\) 等于或不等于 e，g 的阶为 h 或 2h（利用习题 12 f)）。
\item
  元素 \(g \in G\) 是 Coxeter 元，当且仅当 \(Ad g \otimes 1_{C}\)（\(g_{C}\) 的自同构）满足第 VIII 章，§5，习题 5 f) 的诸等价条件。
\end{enumerate}

\(d\) ) 证明每个正则元 \(g\)（属于 \(\mathbf{G}\)，满足 \(g^{h}\in \mathrm{C}(\mathrm{G})\)）都是 Coxeter 元；当 \(p < h\) 时，不存在正则元 \(k\) 满足 \(k^p\in \mathrm{C}(\mathrm{G})\)。

\begin{enumerate}
\def\labelenumi{\arabic{enumi})}
\setcounter{enumi}{14}
\tightlist
\item
  设 H 是 G 的一个连通闭子群。若 H 不包含于任何异于 G 的极大秩连通闭子群之中，则称 H 是净的。
\end{enumerate}

\begin{enumerate}
\def\labelenumi{\alph{enumi})}
\item
  证明 H 是净的，当且仅当它在 G 中的中心化子等于 C(G)。特别地，若 H 是净的，则 C(H) = C(G) ∩ H。
\item
  以下假定 H 是净的。证明对 H 的每个极大环面 S，都有 \(\mathrm{C}(\mathrm{H}) = \mathrm{S} \cap \mathrm{C}(\mathrm{G})\)。
\item
  设 \(H'\) 是 G 的包含 H 的一个连通闭子群。证明 \(rkH < rkH'\)，并由此推出 H 在 \(H'\) 中是净的。
\item
  设 K 是 H 的一个连通闭子群，它在 H 中是净的且包含 G 的一个正则元。证明 K 在 G 中是净的。
\end{enumerate}

\begin{enumerate}
\def\labelenumi{\arabic{enumi})}
\setcounter{enumi}{15}
\tightlist
\item
  设 H 是 G 的一个连通闭子群，使得 \(T \cap H\) 是 H 的一个极大环面 S。设 \(\lambda \in \mathrm{R}(\mathrm{H}, \mathrm{S})\)；用 \(\mathrm{R}(\lambda)\) 表示 \(\mathrm{R}(\mathrm{G}, \mathrm{T})\) 中在 S 上的限制等于 \(\lambda\) 的根组成的集合。
\end{enumerate}

\begin{enumerate}
\def\labelenumi{\alph{enumi})}
\item
  证明 \(\mathrm{R}(\lambda)\) 非空。
\item
  设 \(w \in \mathrm{W}_{\mathrm{H}}(\mathrm{S})\)；证明存在元素 \(\bar{w}\)（属于 \(\mathrm{W}_{\mathrm{G}}(\mathrm{T})\)），使 \(\mathrm{R}(w\lambda) = \bar{w}\mathrm{R}(\lambda)\)（利用 §2 的习题 4）。
\item
  设 \(\mathrm{P}(\lambda)\) 是 \(\mathrm{R}(\mathrm{G},\mathrm{T})\) 与 \(\mathrm{X}(\mathrm{T})\) 中由 \(\mathrm{R}(\lambda)\) 生成的子群之交。证明存在 G 的包含 T 的闭子群 \(G_{\lambda}\)，使 \(\mathrm{P}(\lambda)=\mathrm{R}(\mathrm{G}_{\lambda},\mathrm{T})\)。由此推出反射 \(s_{\lambda}\in\mathrm{W}_{\mathrm{H}}(\mathrm{S})\) 是诸反射 \(s_{\alpha}\)（\(\alpha\in\mathrm{P}(\lambda)\)）之积在 S 上的限制。
\item
  证明结点向量 \(K_{\lambda} \in \mathrm{L}(\mathrm{S})\)（对应于 \(\lambda\)）是诸 \(K_{\alpha}\)（\(\alpha \in \mathrm{R}(\lambda)\)）的一个整系数线性组合。
\item
  设 \(B_{H}\) 是 \(R(H,S)\) 的一个基。证明 \(R(\lambda)\) 包含于 \(X(T)\) 中由诸 \(R(\mu)\)（\(\mu \in B_{H}\)）的并所生成的子群（证明具有所述性质的 \(\lambda \in R(H,S)\) 组成的集合在所有 \(s_{\mu}\)（\(\mu \in B_{H}\)）下稳定，并利用 c)）。
\end{enumerate}

¶ 17) 我们沿用上一习题的记号；再假定子群 H 是净的（习题 15）。

\begin{enumerate}
\def\labelenumi{\alph{enumi})}
\item
  证明 \(\mathrm{R}(\mathrm{G},\mathrm{T})\) 包含于 \(\mathrm{X}(\mathrm{T})\) 中由诸 \(\mathrm{R}(\lambda)\)（\(\lambda\in B_{H}\)）的并所生成的子群。
\item
  设 \(\Delta\) 是 S 中由诸 \(s \in S\)（满足 \(\lambda(s) = \mu(s)\)，对所有 \(\lambda, \mu\) 属于 \(B_{H}\)）组成的子群的单位元连通分支。证明存在 R(G,T) 的一个基 \(\{\alpha_{1}, \ldots, \alpha_{l}\}\) 和一个整数 k，\(0 \leq k \leq l - 1\)，使 \(\Delta\) 是满足下列条件的 \(t \in T\) 组成的集合的单位元连通分支
\end{enumerate}

\[
\alpha_ {1} (t) = \dots = \alpha_ {k} (t) = 1, \alpha_ {k + 1} (t) = \dots = \alpha_ {l} (t).
\]

（设 \(x\) 是 \(\mathrm{L}(\mathrm{S})\) 中使 \(\delta(\lambda)(x) = 2\pi i\)（对所有 \(\lambda \in \mathrm{B}_{\mathrm{H}}\)）成立的元素；由 \(a\)) 推出 \(\exp x \in \mathrm{C}(\mathrm{G})\)。选取 \(\{\alpha_1, \ldots, \alpha_l\}\)，使 \(i\delta(\alpha_j)(x)\) 为零（对 \(1 \leq j \leq k\)）且 \(< 0\)（对 \(k + 1 \leq j \leq l\)）；然后证明对所有 \(\lambda \in \mathrm{B}_{\mathrm{H}}\)，\(\mathrm{R}(\lambda)\) 中的每个根都可写成

\[
\alpha_ {j} + n _ {1} \alpha_ {1} + \dots + n _ {k} \alpha_ {k},
\]

其中 \(j \geq k + 1\) 且 \(n_i \in \mathbf{N}\)。利用 \(a)\) 得出结论。

\begin{enumerate}
\def\labelenumi{\alph{enumi})}
\setcounter{enumi}{2}
\tightlist
\item
  证明整数 k 不依赖于环面 S、T 或基 \(B_{H}, \{\alpha_{1}, \ldots, \alpha_{l}\}\) 的选取；它等于 \(Z_{\mathrm{G}}(\Delta)\) 的换位子群的秩。
\end{enumerate}

\begin{enumerate}
\def\labelenumi{\arabic{enumi})}
\setcounter{enumi}{17}
\tightlist
\item
  我们沿用习题 16 与 17 的记号。若子群 H 是净的，且子环面 \(\Delta\) 包含一个正则元（换言之，k = 0），则称子群 H 是主的。
\end{enumerate}

\begin{enumerate}
\def\labelenumi{\alph{enumi})}
\item
  设 \(\mathrm{K}\) 是 \(\mathrm{H}\) 的一个连通闭子群。证明 \(\mathrm{K}\) 是 \(\mathrm{G}\) 的主子群，当且仅当 \(\mathrm{K}\) 在 \(\mathrm{H}\) 中是主的且 \(\mathrm{H}\) 在 \(\mathrm{G}\) 中是主的（利用习题 15 c) 与 d)）。
\item
  以下假定 H 是主的。证明诸 \(\mathrm{R}(\mu)\)（\(\mu \in B_{H}\)）的并是 \(\mathrm{R}(\mathrm{G}, \mathrm{T})\) 的一个基。
\item
  设 \(\alpha \in \mathrm{R}(\mathrm{G},\mathrm{T})\)。证明 \(\alpha\) 在 S 上的限制是某个根 \(\lambda \in \mathrm{R}(\mathrm{H},\mathrm{S})\) 的非零整数倍；根 \(\alpha\) 是 \(\mathrm{R}(\lambda)\) 中元素之和（证明 \(\mathrm{L}(\mathrm{S})\) 与超平面 \(\delta (\alpha) = 0\) 的交是 \(\mathrm{L}(\mathrm{S}))\) 的一面墙）。
\item
  设 \(\lambda\in\mathrm{R}(\mathrm{H},\mathrm{S})\)；证明结点向量 \(K_{\lambda}\) 是诸 \(K_{\alpha}\)（\(\alpha\in\mathrm{R}(\lambda)\)）的一个非负整系数线性组合（参见~习题 16 d) 以及第 V 章，§3，节 5，引理 6）。
\end{enumerate}

\begin{enumerate}
\def\labelenumi{\arabic{enumi})}
\setcounter{enumi}{18}
\tightlist
\item
  我们沿用习题 16 至 18 的记号；假定子群 H 是主的。给 \(\mathrm{X(T)\otimes\mathbf{R}}\)（相应地 \(\mathrm{X(S)\otimes\mathbf{R}}\)）赋以一个在 \(\mathrm{W_{G}(T)}\)（相应地 \(\mathrm{W_{H}(S)}\)）下不变的内积。设 \(\lambda,\mu\) 是 \(B_{H}\) 中的两个根。
\end{enumerate}

\begin{enumerate}
\def\labelenumi{\alph{enumi})}
\item
  证明若 \(\lambda\) 与 \(\mu\) 正交，则集合 \(\mathrm{R}(\lambda)\) 与 \(\mathrm{R}(\mu)\) 正交。由此推出，若 \(\mathrm{G}\) 是概单的，则 \(\mathrm{H}\) 也是概单的。
\item
  以下假定 \(n(\lambda, \mu) = -1\)。证明存在一个满射 \(u: \mathrm{R}(\lambda) \to \mathrm{R}(\mu)\)，使得对 \(\alpha \in \mathrm{R}(\lambda)\)，\(u(\alpha)\) 是 \(\mathrm{R}(\mu)\) 中与 \(\alpha\) 相连（即不正交）的唯一根；我们有 \(K_{\mu} = \sum_{\beta \in \mathrm{R}(\mu)} K_{\beta}\)，且 \(\mathrm{R}(\mu)\) 中的根两两正交（把 \(K_{\mu}\) 与 \(K_{\lambda}\) 写成诸 \(K_{\alpha}\) 的线性组合，然后比较系数）。
\item
  若 \(n(\mu, \lambda) = -1\)，则映射 \(u\) 是双射，且 \(\mathrm{R}(\lambda) \cup \mathrm{R}(\mu)\) 中相连的根对只有诸对 \((\alpha, u(\alpha))\)（\(\alpha \in \mathrm{R}(\lambda)\)）。
\item
  假定 \(n(\mu,\lambda)=-2\)；设 \(\beta\in\mathrm{R}(\mu)\)。证明 \(u^{-1}(\beta)\) 含有一个或两个元素；若 \(u^{-1}(\beta)=\{\alpha_{1},\alpha_{2}\}\)，则根 \(\alpha_{i}\)（i=1,2）与 \(\mathrm{R}(\lambda)\) 中其余的根正交，且与 \(\beta\) 等长；若 \(u^{-1}(\beta)=\{\alpha\}\) 且 \(\|\alpha\|=\|\beta\|\)，则 \(\alpha\) 与唯一的根 \(\alpha^{\prime}\in\mathrm{R}(\lambda)\) 相连，后者与 \(\alpha\) 等长，且 \(\{\alpha,\alpha^{\prime}\}\) 与 \(\mathrm{R}(\lambda)\) 的其余部分正交；若 \(u^{-1}(\beta)=\{\alpha\}\) 且 \(\|\alpha\|\neq\|\beta\|\)，则 \(\alpha\) 与 \(\mathrm{R}(\lambda)\) 中其余的根正交。
\item
  类似地讨论 \(n(\mu, \lambda) = -3\) 的情形。
\end{enumerate}

\begin{enumerate}
\def\labelenumi{\arabic{enumi})}
\setcounter{enumi}{19}
\tightlist
\item
  假定群 G 是概单的；设 H 是 G 的一个秩 \(\geq 2\) 的主子群。
\end{enumerate}

\begin{enumerate}
\def\labelenumi{\alph{enumi})}
\item
  证明 H 是 \(B_{h}, C_{h}, F_{4}\) 或 \(G_{2}\) 型的半单群（利用习题 19）。
\item
  假定 H 的秩 \(\geq 3\)。证明 G 是 \(A_{l}, D_{l}, E_{6}, E_{7}\) 或 \(E_{8}\) 型的（考察 R(G, T) 的邓金图的端点，并应用习题 19）。
\item
  证明若 \(rkH \geq 3\)，则我们处于下列情形之一：
\end{enumerate}

G 是 \(A_{2l}\) 型 \((l \geq 3)\) 的，H 是 \(B_{l}\) 型的；

G 是 \(\mathrm{A}_{2l - 1}\) 型 \((l\geq 3)\) 的，H 是 \(\mathbf{C}_l\) 型的

G 是 \(D_{l}\) 型 \((l \geq 4)\) 的，H 是 \(B_{l-1}\) 型的；

G 是 \(\mathrm{E}_6\) 型的，H 是 \(\mathrm{F}_4\) 型的。

\(d\) ) 证明若 H 是 \(\mathrm{B}_2\) 型的，则 G 是 \(\mathrm{A}_3\) 或 \(\mathrm{A}_4\) 型的。

\begin{enumerate}
\def\labelenumi{\alph{enumi})}
\setcounter{enumi}{4}
\tightlist
\item
  证明若 H 是 \(G_{2}\) 型的，则 G 是 \(B_{3}, D_{4}\) 或 \(A_{6}\) 型的。
\end{enumerate}

（关于这些情形更明确的描述，参见 §5，习题 5。）

\(f)\) 设 \(\mathrm{K}\) 是 \(\mathrm{G}\) 的包含 \(\mathrm{H}\) 的一个连通闭子群；证明或者 \(\mathrm{K} = \mathrm{G}\)，或者 \(\mathrm{K} = \mathrm{H}\)，或者 \(\mathrm{H}\) 是 \(\mathrm{G}_2\) 型的、\(\mathrm{K}\) 是 \(\mathrm{B}_3\) 型的且 \(\mathrm{G}\) 是 \(\mathrm{D}_4\) 或 \(\mathrm{A}_6\) 型的。

\begin{enumerate}
\def\labelenumi{\arabic{enumi})}
\setcounter{enumi}{20}
\tightlist
\item
  \begin{enumerate}
  \def\labelenumii{\alph{enumii})}
  \tightlist
  \item
    设 H 是 G 的一个秩为 1 的连通闭子群。证明下列条件等价：
  \end{enumerate}
\end{enumerate}

\begin{enumerate}
\def\labelenumi{(\roman{enumi})}
\item
  H 是主的；
\item
  H 是净的且包含 G 的一个正则元；
\item
  存在一个主 \(\mathfrak{sl}_2\)-三元组 \((x,h,y)\)，它位于 \(\mathfrak{g}_{\mathbf{C}}\) 中（第 VIII 章，§11，节 4），使得
\end{enumerate}

\[
\mathrm{L} (\mathrm{H}) _ {(\mathbf {C})} = \mathbf {C} x + \mathbf {C} h + \mathbf {C} y.
\]

\begin{enumerate}
\def\labelenumi{\alph{enumi})}
\setcounter{enumi}{1}
\item
  证明 G 包含一个秩为 1 的主子群，且任意两个这样的子群都共轭（用习题 17 的记号，注意 \(S = \Delta\) ）。
\item
  证明 G 的一个连通闭子群是主的，当且仅当它包含 G 的一个秩为 1 的主子群（利用习题 18 a)）。
\end{enumerate}

\begin{enumerate}
\def\labelenumi{\arabic{enumi})}
\setcounter{enumi}{21}
\item
  设 H 是 G 的一个秩为 1 的主子群；设 \(\Gamma\) 是 \(\operatorname{Aut}(G)\) 中由满足 \(u(\mathrm{H}) = \mathrm{H}\) 的自同构 u 组成的子群。证明 \(\operatorname{Aut}(G) = \Gamma.\operatorname{Int}(G)\)。
\end{enumerate}

